% Put the summary and conclusion here

\section{SUMMARY AND CONCLUSION}

International School Manila faced persistent inconsistencies in sponsorship billing and Letter of Guarantee coverage decisions that stemmed from manual interpretation of LoG provisions at the point of charge entry. Operating across four independent systems—PowerSchool (Student Information System), Student Charging Portal, NetSuite (Enterprise Resource Planning), and the Online Billing System—staff members interpreted similar sponsor coverage agreements differently, leading to billing errors, posting reversals, delayed reconciliation, and audit trail fragmentation. These operational challenges carried compliance risk under the Bureau of Internal Revenue's Ease of Paying Taxes requirements, which demand accurate, well-documented billing records.

This project designed, developed, and deployed a pilot Sponsor Management and Letter of Guarantee Coverage Integration System that centralizes sponsor master data and applies deterministic coverage rules at the point of charge entry. The system was implemented as an ASP.NET Core 8.0 progressive web application hosted on Microsoft Azure with Azure SQL Database, Azure Active Directory authentication (including Google OAuth for staff), and REST APIs documented via Swagger/OpenAPI. Rather than replacing ISM's existing core systems, the solution introduces a shared institutional service that evaluates every charge according to explicit, versioned rules and returns authoritative Bill-To designations and sponsor-parent allocations that propagate to the ERP and statement layer.

\subsection{Achievement of Objectives}

The project successfully achieved all five specific objectives outlined in the Introduction:

\begin{enumerate}
    \item \textbf{Objective 1: Synchronize sponsor management across PowerSchool, NetSuite, the Student Charging Portal, and the Online Billing System.}
    
    This objective was achieved through implementation of a centralized Sponsor Master repository that maintains canonical sponsor records with unique identifiers (SponsorID), sponsor attributes (name, legal name, TIN, address), and contact information. The pilot demonstrated identifier consistency through API endpoints designed for cross-system synchronization. While scheduled integration with PowerSchool, NetSuite, and OBS was deferred to post-pilot production deployment, the API contracts and data structures were validated to support future synchronization workflows driven by business events (create, update, merge).
    
    \item \textbf{Objective 2: Achieve real-time, explainable coverage evaluation based on a deterministic rules engine.}
    
    This objective was fully achieved through implementation of a coverage evaluation API that accepts charge line parameters (student, sponsor, item, amount, date) and returns deterministic coverage decisions (Covered, Not Covered, Split) with human-readable reason codes (e.g., \texttt{NO\_MATCHING\_RULE}, \texttt{COVERAGE\_APPLIES}, \texttt{CAP\_EXCEEDED}) in under 200 milliseconds. The rules engine implements item-level and category-level coverage logic with effective date enforcement, coverage caps, and student-specific overrides. Every evaluation produces an audit record capturing inputs, decision outcome, rule version, and timestamp, ensuring full traceability for billing verification and compliance review. Testing confirmed 100\% consistency: identical inputs always produce identical outputs.
    
    \item \textbf{Objective 3: Create a centralized Sponsor Master repository.}
    
    This objective was successfully implemented with complete CRUD operations for sponsor records, multiple contact management per sponsor, duplicate detection logic by name and TIN, and data persistence across application restarts validated in both local and Azure environments. The Sponsor Master maintains sponsor data, coverage policies (via LoG records), limits (via coverage item definitions), effective dates, status tracking, and cross-system identifiers. The pilot demonstrated that four test accounts (Administrator, Admissions, Cashier, Sponsor) could successfully create, retrieve, update, and search sponsor records, establishing the foundation for PowerSchool and NetSuite identifier propagation in production deployment.
    
    \item \textbf{Objective 4: Ensure correct Bill-To designation is set up in NetSuite and the Online Billing System.}
    
    This objective was partially achieved in the pilot. The coverage evaluation engine correctly computes Bill-To designations (Sponsor, Parent, or Split) and sponsor-parent allocations based on deterministic rules, with results structured for NetSuite posting and statement display. However, actual integration with NetSuite for AR transaction creation and OBS for statement propagation was not implemented in the pilot phase. The API contracts were designed and documented to support these integrations, and the allocation formulas were validated through testing scenarios. Production deployment will complete this objective by implementing scheduled or event-driven synchronization that posts coverage decisions to NetSuite general ledger and presents split allocations in OBS statements with EoPT-compliant documentation.
    
    \item \textbf{Objective 5: Enforce data quality and governance within the Sponsor Master.}
    
    This objective was achieved through implementation of required field validation (SponsorID, name, legal name, TIN), duplicate detection logic that identifies potential matches by name and TIN, role-based access control restricting sponsor creation and editing to Administrator and Admissions roles, and maintenance of audit logs for all sponsor record changes. The system validates input formats, enforces uniqueness constraints on SponsorID, and maintains clean cross-system identifiers suitable for propagation to PowerSchool and NetSuite. Testing confirmed that invalid or incomplete sponsor records are rejected with clear validation messages, and duplicate detection surfaces potential matches for administrative review and merge workflows.
\end{enumerate}

\subsection{Key Contributions}

This work contributes to both academic understanding of integration patterns and operational practice at International School Manila.

\textbf{Academic Contribution:} The project demonstrates a practical integration architecture for multi-system sponsorship management that preserves existing platform investments while resolving inconsistency at its source. Rather than advocating for monolithic ERP replacement or accepting continued reliance on manual spreadsheets, the solution positions a lightweight, purpose-built service as a shared institutional resource. This approach is generalizable to other institutions facing similar multi-system complexity in third-party billing, tuition management, or financial aid administration. The emphasis on deterministic, explainable decision-making aligns with regulatory compliance requirements and offers an alternative to black-box machine learning models in contexts where auditability and traceability are paramount.

\textbf{Institutional Contribution:} International School Manila now has a functional pilot system that validates core sponsor management and coverage evaluation concepts. The pilot provides:

\begin{itemize}
    \item A canonical sponsor master with clean identifiers ready for propagation to downstream systems
    \item A deterministic rules engine capable of real-time coverage evaluation with sub-second response time
    \item Role-based workflows for Administrator, Admissions, Cashier, and Sponsor users
    \item Comprehensive audit trails for all coverage decisions, supporting post-hoc review and dispute resolution
    \item REST APIs documented via Swagger/OpenAPI, enabling future integration with Student Charging Portal and external systems
    \item Security-hardened implementation aligned with OWASP best practices, including HTTPS enforcement, CSRF protection, secure session management, and Google OAuth domain restriction for staff access
\end{itemize}

Expected operational benefits include reduced manual verification effort, fewer posting reversals due to inconsistent Bill-To assignments, faster billing cycles by eliminating back-and-forth clarification between Admissions and Finance, and improved audit readiness through tamper-evident decision logs that satisfy BIR Ease of Paying Taxes documentation requirements.

\textbf{Methodological Contribution:} The project employed a DevOps-inspired development approach with iterative feature delivery, automated testing, security scanning, and controlled deployment to Azure. The progression from 78\% test pass rate in early validation to 86.4\% pass rate in final testing, coupled with identification and remediation of 10 security vulnerabilities via OWASP ZAP scanning, demonstrates the value of continuous testing and security validation throughout development rather than as a late-stage gate. This methodology is transferable to other institutional software development initiatives where quality, security, and maintainability are critical.

\subsection{Contributions to Practice}

The pilot deployment validates that institutions can achieve consistent, auditable sponsorship coverage without replacing core systems. By relocating the decision locus to charge entry and establishing a single source of truth for sponsor data, the solution reduces operational variability and improves financial control. The explicit design for explainability—every decision includes a reason code, rule version, and audit trail—supports staff training, dispute resolution, and regulatory compliance in ways that opaque spreadsheet calculations cannot.

The decision to implement dual authentication (local credentials for initial testing plus Google OAuth for staff) reflects practical institutional requirements. While local accounts support isolated testing and development, Google OAuth with \texttt{@ismanila.org} domain restriction aligns with ISM's existing identity management and simplifies user provisioning. This hybrid approach provides flexibility during pilot validation while establishing a path toward centralized identity management in production.

\subsection{Limitations and Boundaries}

The pilot implementation focused on core sponsor management and coverage evaluation rather than full end-to-end integration. External system synchronization with PowerSchool, NetSuite, and the Online Billing System remains unimplemented, requiring additional development before the solution can drive actual billing workflows. The coverage evaluation engine was tested with synthetic data and small record sets (fewer than 10 sponsors, fewer than 5 students); production-scale performance with hundreds of sponsors and thousands of charges requires validation under realistic load.

The Azure deployment achieved 87.5\% operational status with one known non-critical issue: the coverage preview endpoint returns HTTP 500 errors on Azure production (\url{https://ismsponsor.azurewebsites.net/api/v1/coverage/preview}) due to missing seed data in the Azure SQL database Items and ItemCategories tables. This issue does not impact the pilot demonstration as the local environment provides full functionality with sub-200ms response times. Post-pilot remediation requires executing comprehensive database seeding scripts on Azure SQL before production cutover.

Security testing using OWASP ZAP identified and remediated 10 vulnerability classes across security misconfigurations, cryptographic failures, and access control issues. All vulnerabilities were resolved before final testing, achieving 100\% pass rate on security validation. High availability, disaster recovery, advanced monitoring, and operational runbooks are recommended before full production deployment to ensure system reliability and supportability.

The pilot used in-memory database storage in local development for rapid iteration, with Azure SQL Database serving production. Historical sponsor records and retroactive restatement of past coverage decisions were not included in scope; the system operates as a forward-looking service for new charges and sponsor agreements.

\subsection{Conclusion}

The Sponsor Management and Letter of Guarantee Coverage Integration System successfully addresses the root cause of billing inconsistencies by relocating coverage decision-making to the point of charge entry and establishing a canonical Sponsor Master synchronized across ISM's existing platforms. The pilot implementation demonstrates that institutions can achieve consistent, auditable sponsorship coverage without replacing PowerSchool, NetSuite, or in-house systems, while improving compliance with regulatory requirements such as the Bureau of Internal Revenue's Ease of Paying Taxes framework.

Comprehensive testing validated functional correctness (86.4\% pass rate across 110 test cases), security posture (100\% pass rate on 15 security controls), and performance adequacy (sub-200ms coverage evaluation in local environment). The identification and remediation of 10 security vulnerabilities through OWASP ZAP scanning established a security-first development discipline applicable to future projects.

Expected operational benefits include reduced manual verification, fewer posting reversals, faster billing cycles, and improved audit readiness. These outcomes support ISM's commitment to accuracy, transparency, and accountability in financial operations while preserving the investments in PowerSchool, NetSuite, and in-house systems that serve multiple institutional functions beyond sponsorship billing.

The project establishes a foundation for future enhancements, including production integration with external systems, advanced analytics on sponsor portfolio trends, automated reconciliation workflows, and extended sponsor self-service features. The pilot validates core architectural concepts—centralized sponsor master, deterministic rules engine, API-first integration design, and comprehensive audit trails—that position ISM to scale its sponsorship program while maintaining operational excellence and regulatory compliance.